% Zdroj: PDF priklady-leto-2023.pdf, fyzická str. 1. Značka: neznačená starší sada (léto 2023). Zařazeno: I2.
\section{Třídění sléváním}
\szzotazka{léto 2023}{1}{Třídění sléváním (Mergesort, slévání $k$ posloupností)}

\noindent\textbf{Zadání.}\par
\begin{enumerate}
\item Formulujte algoritmus třídění posloupnosti $n$ celých čísel sléváním (Mergesort)
  a zapište ho v pseudokódu.
\item Analyzujte časovou a prostorovou složitost tohoto algoritmu.
\item Navrhněte efektivní algoritmus pro slévání $k$ setříděných posloupností o~celkem
  $n$ prvcích. Upravte Mergesort, aby toto slévání používal. Jak se změní jeho časová
  složitost? Zapište ji jako funkci proměnných $n$ a $k$.
\end{enumerate}

\medskip
\noindent\textbf{Náčrt řešení.}\par
\textbf{1)} Mergesort (třídění sléváním):
\begin{lstlisting}
MergeSort(a[0..n-1]):
  if n <= 1: return a
  m = n div 2
  L = MergeSort(a[0..m-1]); R = MergeSort(a[m..n-1])
  return Merge(L, R)

Merge(L, R):                 // obe posloupnosti setridene
  i = 0; j = 0; out = []
  while i < |L| and j < |R|:
    if L[i] <= R[j]: out.append(L[i]); i++
    else:            out.append(R[j]); j++
  pripoj zbytek L[i..] a R[j..] na konec out
  return out
\end{lstlisting}

\textbf{2)} \textsc{Merge} přesune každý prvek jednou, stojí $\Theta(n)$. Rekurence
$T(n)=2T(n/2)+\Theta(n)$ má podle kuchařkové věty ($a=2$, $b=2$, $c=1$, $a/b^c=1$)
řešení $T(n)=\Theta(n\log n)$, a~to v~každém případě (nezávisí na datech). Paměť:
pomocná pole pro slévání $\Theta(n)$ a~zásobník rekurze $\Theta(\log n)$, celkem
$\Theta(n)$.

\textbf{3)} \emph{Slévání $k$ posloupností:} do minimové haldy vložíme první prvek
každé posloupnosti (s~informací, odkud pochází). Opakovaně vyjmeme minimum
(\textsc{ExtractMin}), zapíšeme ho na výstup a~vložíme do haldy další prvek téže
posloupnosti, pokud existuje. Halda má nejvýše $k$ prvků, každá operace stojí
$O(\log k)$, každý z~$n$ prvků se jednou vloží a~jednou vyjme: celkem $O(n\log k)$
(plus $O(k)$ na postavení haldy).

\emph{$k$-cestný Mergesort:} vstup rozdělíme na $k$ částí velikosti $n/k$, každou
setřídíme rekurzivně a~výsledky slijeme haldou. Rekurence
$T(n)=k\,T(n/k)+\Theta(n\log k)$. Strom rekurze má $\log_k n$ hladin a~na každé je
slévání celkem $\Theta(n\log k)$, takže
\[
T(n)=\Theta(n\log k\cdot\log_k n)=\Theta\Big(n\log k\cdot\frac{\log n}{\log k}\Big)
=\Theta(n\log n).
\]
Asymptotická složitost se tedy nezmění, faktor $\log k$ se vykrátí (pro $k\ge2$; pro
$k=n$ jde vlastně o~třídění haldou, opět $\Theta(n\log n)$).

\medskip
\textit{(V~PDF bez náčrtu řešení; náčrt doplněn k~výkladu okruhu.)}
